\chapter{The main theorem and the finite group consequence}
\label{chap:current-main}

Manuscript Theorem 1.1 asserts the inductive blockwise Alperin weight
condition for every nonabelian finite simple group of Type B or Type C
and every sporadic simple group, at each prime dividing its order. The
Type B clause includes ranks one and two and groups over fields of even
characteristic. The arguments of the preceding chapters establish the
three cases under their stated hypotheses: manuscript Theorem 3.1,
Corollary 4.15 and Theorem 5.1.

\section{Combining the three families}

Fix an identification of the simple group $S$ with a group in one of
these families. For Type C, use the
specified projective symplectic group over its finite field, with the
exceptional nonsimple case excluded. For Type B, use the specified
orthogonal group in rank at least one. The cases in small ranks and
even characteristic combine with the theorem in higher rank over odd
fields. The latter uses the defining and odd nondefining cases and
the arguments for $\ell=2$. For a sporadic group, use its specified realisation and the cover
attached to the chosen prime.

The theorem for each family gives the full inductive condition on the
specified cover. Compose its covering projection with the chosen
isomorphism to $S$. The covering group, characters, roots and compatible
family then give the condition for $S$. The isomorphism preserves group
order and hence divisibility by $\ell$.

The Type B case thus follows from the hypotheses of the individual
cases and the division into cases by rank and field. The sporadic case combines the
arguments for the four groups treated in Chapter~\ref{chap:current-sporadic}
with the result reported by CTBlocks for the other groups.

\begin{implementation}
\module{ManuscriptIBAW.MainClassical} defines
\lean{Main.ClassicalInputs} and proves
\lean{Main.typeC_complete} and \lean{Main.typeB_complete}.
The latter uses the Type B classification of cases and the theorem
in higher rank obtained from \lean{OddHighInputs.complete}.
\lean{Main.of_allPrimeFamily} uses a \lean{FamilyCertificate}, which
contains the covering group, its projection, an isomorphism from its simple quotient to $S$, and a \lean{FamilyWitness}. The classical conclusion of
\lean{Main.theorem_1_1} requires these full family data. Its sporadic
alternative retains the named group, model and full inductive condition.
\module{ManuscriptIBAW.Main} uses
\lean{Main.Presentation} and \lean{SupportedPresentation}
to prove the three cases in \lean{Main.theorem_1_1}.
\end{implementation}

\section{Reduction to the simple sections of a finite group}

Let $H$ be finite and let $\ell$ be prime. A nonabelian simple group
involved in $H$ is isomorphic to a quotient $U/N$, where $U\leq H$,
$N\lhd U$, and $U/N$ is nonabelian simple. The hypothesis of
manuscript Corollary 1.2 gives an identification with a Type B, Type C or
sporadic simple group for each such quotient whose order is divisible by
$\ell$.

If $\ell$ divides $|U/N|$, apply the main theorem using this identification.
Otherwise, the universal $\ell'$-covering group of $U/N$ has order
prime to $\ell$. Every $\ell$-block of that cover has defect zero,
and the correspondence at $Q=1$ satisfies the inductive condition,
as in the proof of manuscript Corollary 1.2.

Sp\"ath's reduction theorem then gives the blockwise Alperin weight
equality for $H$ \cite[Theorem~A]{Spath2013}:
\[
 |\IBr(b)|=|\Alp(b)|\qquad\text{for every block }b.
\]
Here $\IBr(b)$ consists of the irreducible Brauer characters in $b$,
and $\Alp(b)$ consists of the conjugacy classes of weights belonging
to $b$, with their ordinary local characters of defect zero. The
specified local block operations, character realisations and root
choices describe these objects.

To apply the reduction, each complete family must satisfy the
inductive condition in that theorem. In the exceptional Type B case,
an explicit transport assumption first interprets the blockwise
criteria on the natural quotients as relative block conditions on
the common family over $3.\Omega_7(3)$. This includes character and
block identifications over the fixed coefficient fields, with compatible
root correspondences for quotient groups, extension groups and
intermediate subgroups. A separate published passage combines those
relative block conditions into a full family
\cite[Lemma~3.3]{KoshitaniSpath2016}, with the normalisation of
Sp\"ath's Definition 5.17 and Remark 5.18 \cite{Spath2013}.
This passage chooses compatible global correspondences when necessary.
The proof of Koshitani--Sp\"ath's lemma also supplies the abstract lifting of
characters, weights and blocks from faithful central quotients. Interpreting
these objects in the fixed common family, with the coefficient and root
identifications described above, remains the separate assumption
\lean{NaturalQuotientTransport}.

The formal proof applies the published reduction under these
identifications and the defect zero implication for sections of
order prime to $\ell$. The general reduction theorem is an explicit
external hypothesis. The deduction supplies its hypotheses for every
nonabelian simple section of $H$.

\begin{implementation}
\lean{Main.corollary_1_2} constructs
\lean{SectionEvidence.supported} from the main theorem when
$\ell$ divides the section order and \lean{SectionEvidence.primeTo}
otherwise. It applies
\lean{Main.FiniteGroupReductionSource.reduce}.
The source is indexed by the block family, its character and root
data, and \lean{SpathCoefficientField}. This requires the algebraically
closed modular field $k$ to be algebraic over its prime field $\F_\ell$.
Its antecedent uses the local
\lean{Main.SectionEvidence}, which requires the full classical or
sporadic conclusion on each actual simple section, or that its order is
prime to the chosen prime. The external assumption combines the published reduction with the
defect zero implication for sections of order prime to $\ell$.
Its conclusion is \lean{BlockwiseAWC} for that family.
\end{implementation}
